% Глава: 06-man
% Статус: наполнено (первый прогон).
\chapter{Документация: man, info, help}
\label{ch:06-man}

\section{Зачем это важно}

В Linux основная документация по программам лежит \emph{на самой
системе}: команды сопровождаются справочными страницами, которые
открываются утилитой \cmd{man}. Научиться читать \cmd{man} ---
навык важнее, чем заучивать команды: справочник позволяет
разобраться в любой утилите, даже о которой вы не слышали.

Все справочные страницы доступны прямо на вашей ВМ на
\fpath{pglab.nsu.ru}: \cmd{man} читает документацию, установленную
на сервере.

\begin{tipbox}
Не пытайтесь запомнить все флаги. Запомните \emph{паттерн}:
«откройте \cmd{man}, найдите нужный флаг, прочитайте его
описание». Это покрывает 90\% бытовых задач.
\end{tipbox}

\section{Базовое использование}

\begin{minted}[bgcolor=gray!6]{bash}
man ls              # открыть справочную страницу
man -k "regular expression"   # поиск по кратким описаниям (apropos)
man -S 2 ls         # конкретный раздел
man bash            # документация по оболочке
\end{minted}

Внутри \cmd{man}:
\begin{itemize}
  \item \cmd{q} --- выход;
  \item \cmd{/pattern} --- поиск, \cmd{n} / \cmd{N} --- следующее/
    предыдущее совпадение;
  \item \cmd{g} --- в начало, \cmd{G} --- в конец;
  \item пробел / \cmd{b} --- страница вперёд / назад.
\end{itemize}

Этот интерфейс --- стандартная \textbf{pager} (\cmd{less});
он же используется при чтении больших файлов
(см.\ \nameref{ch:07-text-pipelines}).

\section{Структура справочной страницы}

Страница разбита на разделы. Типовой порядок:
\begin{center}
\begin{tabular}{lp{0.6\linewidth}}
\toprule
\textbf{Раздел} & \textbf{Содержимое} \\
\midrule
NAME       & одно-два предложения: что делает программа \\
SYNOPSIS   & краткий синтаксис, часто с основными флагами \\
DESCRIPTION& подробное описание работы \\
OPTIONS    & все флаги и их смысл (ключевой раздел!) \\
FILES      & файлы, с которыми работает программа \\
EXAMPLES   & готовые примеры использования \\
SEE ALSO   & связанные страницы и утилиты \\
\bottomrule
\end{tabular}
\end{center}

Самый полезный раздел для быстрой работы --- \cmd{OPTIONS}:
он содержит все флаги в алфавитном порядке.

\section{Секции man}

Справочные страницы разделены на \textbf{секции} с цифрами
(от 1 до 9):

\begin{center}
\begin{tabular}{cl}
\toprule
\textbf{Секция} & \textbf{Содержимое} \\
\midrule
1 & Пользовательские команды (\cmd{ls}, \cmd{grep}, \cmd{git}) \\
2 & Системные вызовы ядра (\cmd{open}, \cmd{read}, \cmd{fork}) \\
3 & Библиотечные функции (C: \cmd{malloc}, \cmd{printf}) \\
4 & Файлы устройств (\fpath{/dev/console}) \\
5 & Форматы файлов и конфиги (\fpath{/etc/passwd}) \\
8 & Административные команды (\cmd{mount}, \cmd{useradd}) \\
\bottomrule
\end{tabular}
\end{center}

Если одна команда есть в нескольких секциях, в заголовке
пишут \fpath{ls (1)}, \fpath{ls (8)} и т.п. По умолчанию
открывается первая доступная.

\section{info и встроенный \cmd{--help}}

Помимо \cmd{man} существуют другие источники:
\begin{itemize}
  \item \cmd{--help} (или \cmd{-h}) у большинства программ
    --- краткая справка в терминале, без пагина;
  \item \cmd{info <prog>} --- гиперссылочная документация
    (GNU-инструменты: \cmd{make}, \cmd{gcc}, \cmd{git});
  \item \cmd{apropos <term>} --- поиск по кратким описаниям
    (синоним \cmd{man -k});
  \item \cmd{whatis <prog>} --- только раздел NAME.
\end{itemize}

Правило выбора: для быстрой проверки флага ---
\cmd{--help}; для систематического изучения ---
\cmd{man} или \cmd{info}.

\section{Типичные ошибки}
\begin{itemize}
  \item \emph{No manual entry for X} --- программа не
    установлена или у неё нет страницы; проверьте
    \cmd{which X}.
  \item Не выходит из \cmd{man} --- нажмите \cmd{q}.
  \item Искать весь список флагов в интернете, когда
    \cmd{man} стоит на вашей машине.
  \item Пытаться читать \cmd{man} без понимания структуры:
    сначала \cmd{NAME}, потом \cmd{SYNOPSIS}, и только
    затем \cmd{OPTIONS}.
\end{itemize}

\section{Связь с остальным курсом}

\cmd{man} --- универсальный инструмент, который пригодится
везде: в работе с пакетами
(\nameref{ch:08-packages}), C-компилятором
(\nameref{ch:13-c-compilation}), \cmd{git}
(\nameref{ch:25-git}), Docker (\nameref{ch:19-docker-basics}).
Секция~2 --- системные вызовы, о которых подробнее в главе
о процессах (\nameref{ch:12-processes}).

\section{Задачи для самопроверки}
\begin{enumerate}
  \item Найдите в \cmd{man} флаг \cmd{ls}, который показывает
    размеры в человекочитаемом виде. Как он называется?
  \item Объясните назначение секций 1, 2, 5, 8 и приведите
    по одному примеру каждой.
  \item Как быстро узнать, что делает незнакомая команда,
    не читая всю страницу?
  \item Чем \cmd{--help} отличается от \cmd{man} по
    объёму и структуре?
  \item Найдите в документации \cmd{grep} флаг, который
    выводит номер строки.
\end{enumerate}
